%!TEX root = ../../main.tex

\chapter{Durchführung der Experimente}
\label{ch:durchfuehrung}

Nachdem Kapitel~\ref{ch:implementierung} die technische Umsetzung des Untersuchungssystems beschrieben hat, dokumentiert dieses Kapitel die eigentliche Durchführung der Experimente. Es legt offen, in welchem Ablauf die einzelnen Läufe erfolgten (Abschnitt~\ref{sec:versuchsablauf}), welche Messdaten dabei in welcher Form erhoben wurden (Abschnitt~\ref{sec:messdaten}) und unter welchen Bedingungen die Durchführung stattfand, um die Nachvollziehbarkeit und Reproduzierbarkeit der Ergebnisse sicherzustellen (Abschnitt~\ref{sec:reproduzierbarkeit}). Damit schafft es die Grundlage für die anschließende Auswertung in Kapitel~\ref{ch:ergebnisse}.

\section{Versuchsablauf}
\label{sec:versuchsablauf}

Ausgangspunkt jeder Durchführung ist der in Kapitel~\ref{ch:datensatz} beschriebene synthetische Datensatz, der unverändert für alle Läufe verwendet wird, sowie die zu jeder Aufgabe vorab bestimmte Soll-Lösung. Auf dieser gemeinsamen Grundlage werden die Experimente vollständig automatisiert über den in Abschnitt~\ref{sec:architektur} dargestellten Experiment-Runner ausgeführt, sodass jede Konfiguration unter identischen Bedingungen und ohne manuelle Eingriffe durchläuft. Der Ablauf durchmustert die vollständige Faktormatrix systematisch, indem er für jede Kombination aus Modell, Prompting-Strategie und Aufgabe nacheinander die vorgesehenen Wiederholungen ausführt und das Ergebnis unmittelbar festhält \autocite{boulahiaMulticriteriaEvaluationResearch2024,zhaoSurveyLargeLanguage2026}.

Entsprechend den beiden in Abschnitt~\ref{sec:forschungsdesign} eingeführten Ausführungsmodi gliedert sich die Durchführung in mehrere Teilläufe. Im primären Modus der Code-Generierung wird für jede Konfiguration ein Verarbeitungsskript erzeugt, ausgeführt und bewertet; daraus ergeben sich bei vier Modellen, vier Prompting-Strategien, neun Aufgaben und je drei Wiederholungen insgesamt 432 Einzelläufe. Im ergänzenden Direkt-Modus wird jede der neun Aufgaben von jedem der vier Modelle einmalig unmittelbar auf den eingebetteten Daten gelöst, woraus sich 36 weitere Läufe ergeben. Zusätzlich löst die regelbasierte Vergleichspipeline dieselben neun Aufgaben in je einem deterministischen Durchlauf. Jeder einzelne Lauf wird direkt nach seiner Ausführung als eigenständiger Ergebnisdatensatz gespeichert, sodass der Gesamtablauf jederzeit unterbrochen und fortgesetzt werden kann, ohne dass bereits erhobene Ergebnisse verloren gehen oder überschrieben werden. Die auf den externen Modellzugriff zurückzuführenden, vorübergehenden Störungen werden dabei, wie in Abschnitt~\ref{sec:llm-integration} beschrieben, durch eine automatische Wiederholung transienter Fehlaufrufe abgefangen, während inhaltlich fehlerhafte, aber technisch vollständige Ergebnisse bewusst als gültige Messwerte erhalten bleiben \autocite{zhaoSurveyLargeLanguage2026,hassaniRoleChatGPTData2023}.

\section{Erhobene Messdaten}
\label{sec:messdaten}

Für jeden Lauf wird ein umfassender Ergebnisdatensatz erhoben, der eine spätere Auswertung ohne erneute Modellaufrufe ermöglicht. Er umfasst zum einen die vollständige Konfiguration des Laufs, also das eingesetzte Modell mit seiner exakten Versionskennung, die verwendete Prompting-Strategie, die bearbeitete Aufgabe mit Kategorie und Schwierigkeitsgrad, den zugrunde liegenden Zufallskern, die Nummer der Wiederholung sowie die gewählte Temperatur. Zum anderen enthält er die Ergebnisse der in Abschnitt~\ref{sec:evaluationskriterien} definierten Bewertung, namentlich die Genauigkeit, die Vollständigkeit und die Konsistenz sowie die zugehörigen strukturellen Kennzeichen wie die Übereinstimmung von Schema und Datensatzanzahl, anhand derer die grundsätzliche Vergleichbarkeit eines Ergebnisses beurteilt wird \autocite{ridzuanReviewDataQuality2024,boulahiaMulticriteriaEvaluationResearch2024}.

Ergänzend zu den Qualitätskennzahlen werden für jeden Lauf die Effizienzgrößen aus Abschnitt~\ref{subsec:effizienz} festgehalten, also der Verbrauch an Eingabe- und Ausgabe-Token, die daraus abgeleiteten Nutzungskosten sowie die Antwortzeit des Modells und, im Modus der Code-Generierung, die Dauer und der Erfolg der anschließenden Ausführung \autocite{changEfficientPromptingMethods2024,2026AIIndex}. Um über die reinen Kennzahlen hinaus auch die konkrete Lösung nachvollziehbar zu machen, wird zusätzlich das jeweils erzeugte Artefakt gesichert, im Code-Generierungs-Modus der vollständige erzeugte Quelltext und im Direkt-Modus die unmittelbar zurückgegebene Ergebnistabelle. Sämtliche Läufe werden in einem einheitlichen, maschinenlesbaren Format als jeweils unveränderlicher Datensatz abgelegt, während die erzeugten Ergebnistabellen im typerhaltenden Parquet-Format gespeichert werden. Insgesamt entsteht so aus den 432 Läufen der Code-Generierung, den 36 Läufen des Direkt-Modus und den neun Durchläufen der regelbasierten Pipeline eine geschlossene Datenbasis, auf der die gesamte Auswertung in Kapitel~\ref{ch:ergebnisse} aufsetzt.

\section{Reproduzierbarkeit und Versuchsbedingungen}
\label{sec:reproduzierbarkeit}

Die Reproduzierbarkeit wurde bereits im Untersuchungsdesign als querschnittliche Dimension verankert (Abschnitt~\ref{sec:evaluationskriterien}) und schlägt sich unmittelbar in den Versuchsbedingungen nieder. Um vergleichbare Bedingungen zu schaffen, werden alle steuerbaren Einflussgrößen über sämtliche Läufe hinweg konstant gehalten: Der Datensatz wird über einen festen Zufallskern reproduzierbar erzeugt und unverändert verwendet, und die Temperatur wird durchgehend auf einen niedrigen Wert festgelegt, um die nicht durch das Modell selbst bedingte Variabilität zu minimieren \autocite{schulhoffPromptReportSystematic2025,zhaoSurveyLargeLanguage2026}. Da sich das Feld der Modelle mit hoher Geschwindigkeit weiterentwickelt und die genaue Modellfassung bei den kommerziell betriebenen Anbietern nicht selbst kontrolliert werden kann, wird die exakte Versionskennung jedes Modells bei jedem Aufruf protokolliert; im Rahmen dieser Arbeit kamen je ein aktueller Vertreter von Anthropic, OpenAI und Google sowie ein lokal ausgeführtes Modell der Llama-Reihe zum Einsatz \autocite{2026AIIndex,zhaoSurveyLargeLanguage2026}.

Um die probabilistische Variabilität der Modelle nicht nur zu kontrollieren, sondern selbst messbar zu machen, wird im primären Ausführungsmodus jede Konfiguration dreifach ausgeführt, sodass sich die Streuung der Ergebnisse über die Wiederholungen als Maß für die Stabilität auswerten lässt und damit die zur Reproduzierbarkeit formulierte Hypothese~H3 überprüfbar wird \autocite{zhaoSurveyLargeLanguage2026,hassaniRoleChatGPTData2023}. Der ergänzende Direkt-Modus wird demgegenüber bewusst nur einfach ausgeführt, um den durch die eingebetteten Daten deutlich erhöhten Token-Verbrauch zu begrenzen; die dort erhobenen Werte sind entsprechend als einzelne Beobachtung und nicht als über Wiederholungen gemittelte Größe zu verstehen. Zu berücksichtigen ist ferner, dass einige aktuelle Modelle bestimmte Sampling-Parameter nicht mehr akzeptieren oder standardmäßig einen internen Denkprozess ausführen; solche anbieterspezifischen Abweichungen werden, wie in Abschnitt~\ref{sec:llm-integration} beschrieben, abgefangen und im Ergebnisdatensatz vermerkt, sodass sie bei der Interpretation sichtbar bleiben. Das lokal betriebene Modell wurde auf allgemein verfügbarer Endgeräte-Hardware mit begrenztem Grafikspeicher ausgeführt, was insbesondere im Direkt-Modus eine Begrenzung des nutzbaren Kontextfensters erforderte und die praktischen Grenzen eines lokalen Betriebs sichtbar macht. Trotz dieser Vorkehrungen bleibt eine vollständige, bitgenaue Reproduktion aufgrund der probabilistischen Natur der Modelle und möglicher serverseitiger Änderungen der kommerziellen Modelle nicht garantierbar; die protokollierten Versionskennungen, die konstant gehaltenen Bedingungen und die vollständige Ablage aller Läufe stellen jedoch eine im Rahmen des Möglichen weitgehende Nachvollziehbarkeit sicher \autocite{zhaoSurveyLargeLanguage2026,hassaniRoleChatGPTData2023}.
